\documentclass[11pt]{article}
\usepackage[margin=1in]{geometry}
\usepackage[T1]{fontenc}
\usepackage{lmodern,amsmath,amssymb,booktabs,graphicx}
\usepackage[colorlinks=true,linkcolor=blue,urlcolor=blue]{hyperref}
\setlength{\parindent}{0pt}\setlength{\parskip}{5pt}
\title{RoCE-K: conditional nuisance control and source sample reuse}
\author{}\date{3 October 2026}
\begin{document}\maketitle

\textbf{Main finding.} Conditional outcome concentration and reuse of source
design observations reduce the protected interval length. Independent
confirmation shows gains over the equally improved two-fold target AIPW
certificate. However, a natural unsplit stratified target-only procedure is
still shorter in most of these finite-stratum experiments. We therefore do
not claim a general practical borrowing advantage.

\section{The refinement and its assumptions}

Keep $\tau=d+r_1-r_0$ and the direct-dispersion source procedure. The true
valid source sets are fixed population properties, although unknown to the
procedure. The coverage proof needs a guarantee for these sets, not for sets
chosen after observing fitted errors.

In the four-stratum diagnostic, propensity and joint weights depend only on
$(X,A)$. Conditional on the designs, target outcome successes in each cell
are independent binomials. For coefficients $c_x$ fixed given these designs,
the fitting drift is $\sum_xc_x\{m_t^a(x)-\hat m^a(x)\}$. A bound
$|c_x|\le C_x$ yields
\[
 \left|\sum_xc_x\{m_t^a(x)-\hat m^a(x)\}\right|
 \le \frac12\sum_x\frac{C_x}{N_{xa}+1}
 +\sqrt{\frac12\log(2/\delta)
             \sum_x\frac{N_{xa}C_x^2}{(N_{xa}+1)^2}}
\]
with the allocated conditional probability. This treats the complete weighted
outcome error, rather than a worst-case outcome box in every coordinate.
It is not valid without further work when weights use these same outcomes.

The principal coefficient certificate uses source design means and dispersions
and does not assume a numerical lower bound for source cell probabilities.
The supplementary MGF method uses the exact half-count binomial training law
and assumes $p_j(x,A=a)\ge p_*$. We prespecify $p_*\in\{.02,.04,.06\}$.
The DGP has minimum source joint probability about .06201; this fact verifies
the simulated assumptions and is not supplied to the estimator. The MGF
bounds use stochastic bracketing over probability intervals, not an unchecked
probability grid. The prototype's computed bounds include a floating-point
cushion; it is not an interval-arithmetic certificate.

Direct dispersion does not choose evaluation-dependent Fourier frequencies
or use pilot variance certificates. Its evaluation and nuisance-certificate
events can therefore be combined by a union bound even when observations
are reused. We compare 500 training + 250 pilot + 250 evaluation against
500 training + 500 evaluation, using those same 500 held-out designs for the
nuisance certificate. Every site still has 1,000 distinct patients. Initial
fits are unchanged. The Fourier implementation explicitly rejects this reuse
option pending a separate proof.

\section{Prespecified comparisons and patient budgets}

The main panel has $K=8,32,128,512$, all-valid or weak-bias sources, and target
heterogeneity 0 or .3: 16 settings with 300 repetitions each. The evaluation
reuse experiment is paired on the same 4,800 datasets. Four independent-seed
confirmation settings, $K=128,512$ with weak or strong bias and heterogeneity
.3, have 1,000 repetitions each. A subsequent exploratory extension has
$K=2048$, all-valid, weak or strong bias, with 300 repetitions each. There are
9,700 distinct setting-specific repetitions; reused datasets are counted once.

The DGP and half-count nuisance fits are retained from the prior report.
One quarter of sources have a treated-arm departure; control arms remain
valid. The weak shift is
$2K^{-1/4}\sqrt{.5/250}$ in both members of each paired comparison, so reuse
does not change the underlying population. The strong shift is .2. Each arm
is assumed to have at least $.75K$ valid sources. The target truth is .10.

All source methods use total failure allowance .05: target .02, nuisance
.005, a reserved variance allowance .005, and .01 per arm for direct
source mean/dispersion inference. The reserved variance allowance is unused
by direct dispersion but is not reassigned after seeing the results.
Methods and probability-lower-bound sensitivities are reported separately;
there is no data-dependent choice of the shortest procedure.

\section{Stronger target-only comparators change the interpretation}

First, the same conditional nuisance improvement is applied to both folds of
the full-target AIPW reference. Second, because the model has only four strata,
we add the unsplit full-target plug-in estimator of cell outcome contrasts.
Its finite-sample interval adds a conditional weighted-outcome Hoeffding bound,
including smoothing bias, and a population-composition bound for true cell
contrasts in $[-1,1]$. It uses all 1,000 target patients without a fitting split.
It is not a normal-approximation comparator and does not use oracle variances.

In the independent confirmation data the improved two-fold target interval
has mean length about .697, whereas the unsplit stratified interval has length
about .395. The table shows the complete source comparison for the two central
candidates. ``Design'' requires no supplied source probability minimum;
``MGF (.04)'' assumes that lower bound. Rows with $K=128,512$ use 1,000
repetitions; $K=2048$ rows are the exploratory 300-repeat extension.

\begin{center}\small
\input{confirmation_table}
\end{center}

At $K=512$ and weak bias, the conditional-design interval with reuse has
mean length .5253, versus .697 for two-fold target AIPW: a reduction near
25\%. But .5253 remains larger than the stratified target length .3951.
The MGF option reaches .4727 and also remains longer. At $K=2048$ with weak
bias, the design method has length .4422 and MGF (.04) has .3927. The latter
is only near the stratified reference and requires additional source overlap
information; it is not evidence of a broad advantage.

\begin{figure}[ht]\centering
\includegraphics[width=\linewidth]{training_comparison.pdf}
\caption{Left: paired weak-bias comparisons at target heterogeneity .3,
including both finite-sample target references. Right: independent-confirmation
midpoint RMSE for the conditional-design method with reuse. The two panels
use different prespecified seeds.}
\end{figure}

\section{Coverage and point-estimation boundaries}

All 234 method-setting rows have observed coverage 1. The valid-subset and
target nuisance checks have no observed failures. These are conservative
results, not evidence of exact nominal coverage. The implementation passes
33 focused tests, including finite-distribution checks of the scalar outcome
bound and MGF envelope, grid bracketing, pooled-patient moment identities,
empty target arms, and rejection of unproved reuse with Fourier.

Weak-bias midpoint RMSE improves: at $K=512$ it is about .0148 versus .0315
for the two-fold target point. Under strong bias, however, midpoint RMSE is
about .0505 and bias about .0477. At $K=2048$ strong-bias midpoint bias is
about .0501. Covering the true effect with a protected interval does not
make its midpoint a robust point estimator. All such cases remain reported.

\section{What remains}

The refinement validates two concrete ideas: exploit the complete conditional
outcome error, and avoid unnecessary separation of design calibration from
direct-dispersion evaluation. It does not yet supply high-dimensional theory
for outcome-dependent calibrated weights. A fixed stratum count and a known
source probability lower bound must not be silently promoted to that setting.

The next derivation should control the complete TATE direction, retaining
cross-arm covariance and shared-training cancellations. Adding separate arm
noise and bias radii can lose information that is available in the contrast.
Any such simplification must preserve arm-specific valid sets and be compared
with both target references. The strong-bias behavior of the midpoint remains
a separate limitation; interval coverage is not a point-estimation guarantee.
\end{document}
